\section{实验讨论}

本节把 NPQR 的工程实现回扣到算法分析和量子信息基础知识。报告的重点不是堆叠公式，
而是说明为什么某些阶段耗时高、为什么 SWAP 数和线路深度有物理意义，以及本项目与
课程内容之间的对应关系。

\subsection{复杂度}

令 $n=|V|$ 表示物理量子比特数，$m$ 表示线路门数，$e=|E|$ 表示硬件边数，$K$
表示初始映射候选数，$B$ 表示束宽，$D$ 表示主搜索最大展开深度，$A$ 表示每个状态
保留的候选动作数，$R$ 表示后缀修复深度。NPQR 的主要成本来自候选映射、动作生成
与掩码、模型推理、束搜索扩展和局部修复。

\begin{table}[H]
  \centering
  \caption{复杂度}
  \label{tab:complexity}
  \begin{tabularx}{\linewidth}{lYY}
    \toprule
    阶段 & 时间复杂度估计 & 说明 \\
    \midrule
    拓扑距离矩阵 & $O(n^3)$ 或 $O(n(e+n\log n))$ & 取决于 Floyd 或多源最短路；固定拓扑下可复用。\\
    线路依赖图构建 & $O(m)$ & 扫描门序列，维护前沿门和执行依赖。\\
    逻辑交互图构建 & $O(m)$ & 统计双比特门交互强度，服务初始映射。\\
    初始映射候选 & $O(K\cdot n^2)$ 量级 & 与候选数量和映射评分方式有关。\\
    动作生成与掩码 & $O(B\cdot e)$ & 每个保留状态在硬件边上筛选候选 SWAP。\\
    模型推理 & $O(B\cdot A\cdot C_\theta)$ & $C_\theta$ 为单次模型评分成本。\\
    束搜索扩展 & $O(D\cdot B\cdot A\cdot C_s)$ & $C_s$ 包括状态复制、执行门和评分更新。\\
    后缀修复 & $O(A^R)$ 的受限搜索 & 只在接近完成时启动，避免全局穷举。\\
    轨迹复放 & $O(m+N_{\mathrm{swap}})$ & 线性验证输出轨迹是否合法。\\
    \bottomrule
  \end{tabularx}
\end{table}


表 \ref{tab:complexity} 中的估计解释了实验计时结果。拓扑距离矩阵和线路依赖图构建
属于固定预处理，规模可控；候选映射和主搜索会随线路结构快速变重。束宽 $B$ 和
动作数 $A$ 增大时，搜索质量可能提高，但状态数量和模型评分次数也随之增加。后缀
修复在最坏情况下呈指数增长，因此只用于接近完成的局部尾部。

\subsection{方法对比}

\begin{table}[H]
  \centering
  \caption{方法对比}
  \label{tab:algorithm-comparison}
  \begin{tabularx}{\linewidth}{lYYY}
    \toprule
    方法 & 核心思想 & 优点 & 本项目中的角色 \\
    \midrule
    SABRE & 基于前沿门和后续门距离选择 SWAP & 稳定、快速、工程常用 & 固定质量基线，不参与 NPQR 路线生成。\\
    QRoute & 图神经网络辅助蒙特卡洛树搜索 & 结合学习模型与树搜索 & 相关方法参考。\\
    AlphaRouter & 强化学习结合搜索策略 & 强调策略学习和多步规划 & 相关方法参考。\\
    NPQR & 模型评分、初始映射、束搜索、剪枝和修复组合 & 可解释、可接口调用、可复放验证 & 项目自研默认路由流程。\\
    \bottomrule
  \end{tabularx}
\end{table}


SABRE 的核心是基于前沿门和后续门距离的启发式 SWAP 选择，优点是速度快、实现成熟。
NPQR 的区别在于把动作评分的一部分交给模型学习，同时保留显式搜索、剪枝、修复和
复放验证。两者的比较说明一个工程事实：神经网络评分需要与拓扑约束、动作掩码和
搜索过程结合，才能产生可复放的合法线路。

\subsection{结果解释}

基于实验和复杂度分析，NPQR 在 10/20 比特代表样例和选定 30/50 比特网格样例上
取得了低于 SABRE basic 的 SWAP 数。它的运行时间主要消耗在候选映射、模型推理和
束搜索上。这个结果与项目定位一致：系统优先展示可复放路线、可解释搜索过程和
多种接口交互能力，并把编译速度优化留给后续工程迭代。
